\documentclass[11pt,a4paper]{article}

\usepackage[margin=2.5cm]{geometry}
\usepackage{booktabs}
\usepackage{longtable}
\usepackage{multirow}
\usepackage{array}
\usepackage{xcolor}
\usepackage{hyperref}
\usepackage{graphicx}
\usepackage{caption}
\usepackage{subcaption}
\usepackage{enumitem}
\usepackage{microtype}
\usepackage{lmodern}
\usepackage[T1]{fontenc}
\usepackage[utf8]{inputenc}
\usepackage{amsmath}
\usepackage{siunitx}
\usepackage{tabularx}
\usepackage{makecell}
\usepackage{pdflscape}
\usepackage{float}
\usepackage[numbers,square,sort&compress]{natbib}

\hypersetup{
  colorlinks=true,
  linkcolor=blue!60!black,
  urlcolor=blue!60!black,
  citecolor=blue!60!black
}

\definecolor{lightgray}{gray}{0.92}
\newcolumntype{R}[1]{>{\raggedleft\arraybackslash}p{#1}}
\newcolumntype{L}[1]{>{\raggedright\arraybackslash}p{#1}}
\newcolumntype{C}[1]{>{\centering\arraybackslash}p{#1}}

\title{\textbf{Structured Data on News \& Fact-Checking Domains:\\
  Extraction and Quality Analysis of WDC Schema.org Annotations}\\[0.5em]
  \large Technical Report}
\author{Marwah Sulaiman}
\date{May 2026}

\begin{document}
\maketitle
\tableofcontents
\newpage

% =============================================================================
\section{Motivation and Scope}
\label{sec:motivation}
% =============================================================================

Structured semantic annotations embedded in web pages allow publishers to
represent page content as machine-readable triples
$\langle\text{subject, predicate, object}\rangle$ using shared vocabularies
such as Schema.org~\citep{guha2016schema}.  Over 52\% of pages in the Common
Crawl web archive already carry such annotations~\citep{commoncrawl2024},
yet their quality remains limited: most Schema.org properties are left
unpopulated, and annotations tend to describe page metadata rather than
actual content.

This work examines that gap specifically for news outlets and fact-checking
organisations, sources that are particularly relevant for automated
fact-checking.  The analysis is intended to motivate and ground a broader
research on using large language models (LLMs) to automatically
extract richer, content-grounded annotations from web pages, whose output
could serve as structured evidence for fact-checking systems.

In this work, we assess the completeness, syntactic correctness, structural complexity, and internal consistency of Schema.org annotations across a curated set of news and fact-checking domains, and conduct a dedicated analysis of the \texttt{ClaimReview}~\citep{schemaorg_claimreview} class used by fact-checking organisations.

The \textbf{Web Data Commons (WDC)} project~\citep{meusel2014wdc} (\url{https://webdatacommons.org})
extracts all structured data embedded in the Common Crawl web archive~\citep{commoncrawl2024} and
publishes it as N-Quad files.  The 2024-12 release is the latest snapshot and
covers the entire crawl; it also provides 50 Schema.org class-specific
subsets, each includes all quads whose source pages declare at least one
entity of that class.

This report documents the full pipeline used to retrieve, process, and
analyse WDC data restricted to a curated list of news and fact-checking
domains, spanning three Jupyter notebooks:

\begin{itemize}[leftmargin=1.5em]
  \item \texttt{wdc\_lookup\_extractor.ipynb}: domain list assembly, chunk
        planning, and per-class extraction;
  \item \texttt{wdc\_analysis\_merged.ipynb}: corpus-wide structural and
        quality analysis via DuckDB;
  \item \texttt{wdc\_analysis\_by\_class.ipynb}: deep dive into the
        \texttt{ClaimReview} class.
\end{itemize}

% % =============================================================================
% \section{Target Domain Assembly}
% \label{sec:domains}
% % =============================================================================

% \subsection{Why a curated domain list?}

% Processing the full WDC corpus to find news-relevant data is impractical for
% three reasons.  First, the WDC servers enforce an IP-level rate limit, making
% parallel downloads from a single machine impossible.  Second, streaming and
% filtering 14+ terabytes of compressed N-Quads is not standard practice in
% research settings and would require substantial cluster infrastructure.  Third,
% preliminary inspection of the WDC per-class lookup tables confirmed that
% news and fact-checking pay-level domains (PLDs) are sparse across the 50 class
% subsets, meaning most of the download effort would yield no relevant data.

% A domain-targeted approach (identify the relevant domains first, then
% download only the chunks containing their data) is therefore both practical
% and more principled.



% =============================================================================
\section{WDC Data Retrieval Strategy}
\label{sec:retrieval}
% =============================================================================
\subsection{Sources}

Targetted domains were assembled from two primary sources:

\paragraph{Fact-checking domains.}
The \textbf{Fact-Check Insights} dataset~\citep{factcheckinsights2024}
(\url{https://factcheckinsights.org/}) is a compilation of all fact-checks
published by organisations verified by the Duke Reporters' Lab that carry
\texttt{ClaimReview} and/or \texttt{MediaReview} markup.  From the
\texttt{author.url} field of \texttt{claim\_review.csv}, the top-15
fact-checking organisation domains ranked by publication frequency were
extracted using \texttt{SourceDomainsExtraction/Fact\_Check/extract\_domains.py}.

\paragraph{News domains.}
The \textbf{Media Cloud Directory API}~\citep{mediacloud2024} (\url{https://www.mediacloud.org})
indexes over 60{,}000 news sources worldwide.  National collections for each
country were filtered to the top 5th percentile by total story count and
stories per week, yielding high-volume national sources from every country.
The output is \texttt{national\_sources\_top5pct\_both\_20260424.csv}, generated
by \texttt{SourceDomainsExtraction/News/mediacloud\_national\_high\_volume.ipynb}.

\vspace{1em}


The two collections were merged into a final list of \textbf{388 domains}, which forms the basis for the retrieval pipeline described in the next subsections.


\subsection{The PLD Lookup Table Approach}

For each of the 50 WDC class-specific subsets, WDC publishes a \textbf{PLD
lookup CSV} that maps pay-level domains to the chunk file numbers that contain
their data.  This makes it possible to have a proper sample without downloading the full corpus: for
each class, only the chunks that include quads of the target domain list
need to be retrieved.

\begin{enumerate}[leftmargin=1.5em]
  \item Download the lookup CSV for every class.
  \item Intersect each lookup with the 388 target domains.
  \item Record which chunks are needed for which class.
  \item Download and filter only those chunks.
\end{enumerate}


Out of 50 classes, \textbf{37 had at least one matching domain}; 13 had none.


\subsection{Priority-Based Chunk Reduction}

The naive plan required downloading \textbf{1{,}093 chunks} across 38 classes,
estimated at roughly one week of sequential downloads.  A priority filter
reduced the scope to \textbf{134 chunks} (37 classes) by applying the following
rules:

\begin{itemize}[leftmargin=1.5em]
  \item At most \textbf{5 chunks per class} are selected.
  \item Chunks are ranked by a score combining the count of matching target
        PLDs they contain and their type: core PLDs (fact-checkers and major
        news outlets) contribute a higher weight to the score than general
        news domains, so chunks covering core sources are preferred.
        Higher-scoring chunks are selected first.
  \item \textbf{Core PLDs} (fact-checkers and major news outlets) bypass the
        5-chunk cap: any chunk containing a core PLD is always included.
  \item Each core PLD is allowed to appear in at most 5 classes to prevent
        excessive duplication across the merged output.
\end{itemize}

The result is an \textbf{88\% reduction} in download volume (1{,}093~$\to$~134
chunks) with full representation of the most important domains in each class.

\subsection{Download Pipeline}

The download script (\texttt{wdc\_domain\_filter.py}) implements several mechanisms suited to a rate-limited environment:

\begin{itemize}[leftmargin=1.5em]
  \item \textbf{Sequential downloads} with a 30-second pause between files.
  \item \textbf{Exponential backoff} starting at 120 seconds on HTTP~429
        responses, with each retry waiting longer than the last.
  \item \textbf{Pipelined execution}: chunk $N{+}1$ is downloaded concurrently
        while chunk $N$ is being filtered, hiding download latency.
  \item \textbf{Checkpointing} to \texttt{checkpoint\_lookup.json}: the pipeline
        can be safely interrupted and resumed.
  \item \textbf{Immediate deletion} of each chunk file after filtering to
        minimise disk usage.
\end{itemize}

For each downloaded chunk (a gzip-compressed N-Quads file), every line is
scanned for target domain patterns (\texttt{://domain/}, \texttt{://domain>},
\texttt{.domain/}, \texttt{.domain>}).  Matching lines are appended to
\texttt{output/<ClassName>.nq}.

% =============================================================================
\section{Extraction Results}
\label{sec:results_extract}
% =============================================================================

\subsection{Per-Class Output}

Table~\ref{tab:class_results} summarises the extraction statistics for all 37
classes that yielded data.  Total raw matches (before global deduplication)
sum to \textbf{78{,}940{,}582 quads}.

\begin{longtable}{L{5.6cm} R{2.1cm} R{2.0cm} R{1.5cm}}
\caption{Per-class extraction results. Quads = lines in the per-class
\texttt{.nq} file (before cross-class deduplication).}
\label{tab:class_results}\\
\toprule
\textbf{Class} & \textbf{Quads} & \textbf{File (MB)} & \textbf{Chunks} \\
\midrule
\endfirsthead
\multicolumn{4}{c}{\tablename~\ref{tab:class_results} (continued)}\\
\toprule
\textbf{Class} & \textbf{Quads} & \textbf{File (MB)} & \textbf{Chunks} \\
\midrule
\endhead
\bottomrule
\endfoot
\bottomrule
\endlastfoot
Country                         & 26{,}683{,}479 & 4{,}409.1 & 5 \\
SportsEvent                     &  9{,}633{,}393 & 1{,}754.4 & 4 \\
Place                           &  7{,}807{,}918 & 1{,}580.6 & 4 \\
AdministrativeArea              &  6{,}752{,}407 & 1{,}613.8 & 4 \\
SportsTeam                      &  5{,}786{,}149 & 1{,}019.1 & 4 \\
Event                           &  5{,}380{,}763 & 1{,}163.7 & 5 \\
GeoCoordinates                  &  4{,}070{,}109 &   971.3   & 5 \\
LocalBusiness                   &  3{,}489{,}727 &   609.8   & 5 \\
Person                          &  2{,}433{,}147 &   621.6   & 7 \\
Organization                    &  2{,}150{,}703 &   416.4   & 6 \\
Product                         &  1{,}763{,}672 &   502.9   & 4 \\
QAPage                          &  1{,}313{,}529 &   337.0   & 5 \\
ClaimReview                     &    470{,}994   &   115.4   & 1 \\
Language                        &    372{,}426   &    91.1   & 5 \\
SearchAction                    &    294{,}263   &    73.2   & 5 \\
Answer                          &    265{,}441   &    73.7   & 5 \\
Question                        &    237{,}012   &    68.7   & 5 \\
FAQPage                         &    218{,}073   &    65.4   & 5 \\
Hospital                        &    243{,}676   &    34.5   & 1 \\
Recipe                          &    207{,}258   &    42.6   & 5 \\
JobPosting                      &    719{,}570   &   160.6   & 5 \\
Dataset                         &    141{,}510   &    22.9   & 1 \\
Movie                           &    143{,}037   &    26.9   & 5 \\
StadiumOrArena                  &    736{,}013   &   128.4   & 2 \\
City                            &    763{,}941   &   141.0   & 3 \\
TVEpisode                       &    111{,}867   &    16.7   & 2 \\
RadioStation                    &     67{,}173   &    13.5   & 1 \\
Restaurant                      &     48{,}786   &    11.4   & 5 \\
Book                            &     30{,}572   &     7.6   & 5 \\
MusicRecording                  &     19{,}070   &     3.7   & 1 \\
Hotel                           &      4{,}281   &     0.8   & 1 \\
ShoppingCenter                  &      3{,}599   &     0.7   & 1 \\
CollegeOrUniversity             &      2{,}198   &     0.6   & 2 \\
GovernmentOrganization          &      1{,}517   &     0.3   & 1 \\
LandmarksOrHistoricalBuildings  &        624     &     0.2   & 1 \\
CreativeWork                    &     87{,}243   &    18.2   & 5 \\
EducationalOrganization         &        252     &     0.06  & 2 \\
\midrule
\textbf{Total}                  & \textbf{78{,}940{,}582} & \textbf{16{,}505} & \textbf{134} \\
\end{longtable}

\subsection{Merging and Deduplication}

All 37 per-class \texttt{.nq} files were merged into
\texttt{output/merged.nq} with \textbf{global exact-line deduplication} across
classes.  

\begin{table}[H]
\centering
\caption{Merged corpus summary.}
\label{tab:merged_summary}
\begin{tabular}{lr}
\toprule
\textbf{Metric} & \textbf{Value} \\
\midrule
Unique quads (lines in \texttt{merged.nq}) & 69{,}621{,}246 \\
Uncompressed file size & 14{,}360 MB (14.4 GB) \\
Unique registered domains & 119 \\
Unique subdomains (hosts) & 363 \\
Distinct Schema.org types declared & 162 \\
\bottomrule
\end{tabular}
\end{table}

\subsection{Domain Coverage}

Table~\ref{tab:domain_coverage} lists the top 15 registered domains in the merged corpus by quad count,
combining page coverage, annotation density, and entity statistics for each outlet.

\begin{table}[H]
\centering
\small
\setlength{\tabcolsep}{4pt}
\caption{Top 15 registered domains by quad count (merged corpus).
  \textit{Pages}: distinct crawled pages;
  \textit{Quads}: total RDF quads extracted;
  \textit{\%Q}: share of the 69.6\,M corpus quads;
  \textit{Q/p}: quads per page (annotation density);
  \textit{Entities}: distinct entities across the domain's pages;
  \textit{\%E}: share of the ${\sim}$12.2\,M corpus entities.}

\label{tab:domain_coverage}
\begin{tabular}{lrrrrrr}
\toprule
\textbf{Domain} & \textbf{Pages} & \textbf{Quads} & \textbf{\%Q} & \textbf{Q/p} & \textbf{Entities} & \textbf{\%E} \\
\midrule
ria.ru              & 101{,}258 & 26{,}812{,}929 & 38.5 & 264.8 & 3{,}638{,}349 & 29.8 \\
elpais.com          &  36{,}821 &  7{,}975{,}385 & 11.5 & 216.6 & 1{,}301{,}395 & 10.7 \\
cbssports.com       &  17{,}983 &  5{,}117{,}711 &  7.4 & 284.6 & 1{,}221{,}854 & 10.0 \\
sports.ru           &  38{,}174 &  3{,}043{,}424 &  4.4 &  79.7 &   723{,}710   &  5.9 \\
freepressjournal.in &  15{,}986 &  2{,}800{,}707 &  4.0 & 175.2 &   675{,}231   &  5.5 \\
indiatimes.com      &  33{,}332 &  2{,}524{,}858 &  3.6 &  75.7 &   503{,}533   &  4.1 \\
rbc.ua              &  30{,}542 &  2{,}455{,}246 &  3.5 &  80.4 &   408{,}279   &  3.3 \\
cnn.com             &  18{,}825 &  1{,}915{,}302 &  2.8 & 101.7 &   281{,}225   &  2.3 \\
cbsnews.com         &  12{,}334 &  1{,}700{,}837 &  2.4 & 137.9 &   251{,}695   &  2.1 \\
lefigaro.fr         &  24{,}615 &  1{,}533{,}654 &  2.2 &  62.3 &   342{,}138   &  2.8 \\
as.com              &  17{,}678 &  1{,}505{,}213 &  2.2 &  85.1 &   340{,}609   &  2.8 \\
elespanol.com       &  13{,}306 &  1{,}284{,}796 &  1.8 &  96.6 &   248{,}921   &  2.0 \\
aif.ru              &  30{,}720 &  1{,}241{,}772 &  1.8 &  40.4 &   266{,}521   &  2.2 \\
wikipedia.org       &  66{,}735 &  1{,}228{,}468 &  1.8 &  18.4 &   277{,}768   &  2.3 \\
dw.com              &  41{,}245 &    826{,}324   &  1.2 &  20.0 &   208{,}339   &  1.7 \\
\midrule
\textit{Remaining 104 domains} & \multicolumn{5}{l}{\textit{each ${<}1.2\%$ of total quads}} & \\
\bottomrule
\end{tabular}
\end{table}

The appendix provides the complete list of domains. It is worth noting that some news sources block the Common Crawl bot from indexing their pages, so not all domains are expected to appear in the retrieved data. Nevertheless, many major and high-value sources are well-represented.

\subsection{Schema.org Type Distribution}

The merged corpus contains 162 distinct Schema.org types.  The top 20 by
entity count are shown in Table~\ref{tab:types_merged}.

\begin{table}[H]
\centering
\caption{Top 20 declared entity types in the merged corpus.}
\label{tab:types_merged}
\begin{tabular}{lr}
\toprule
\textbf{Type} & \textbf{Entities} \\
\midrule
ImageObject              & 1{,}746{,}687 \\
ListItem                 & 1{,}509{,}100 \\
SportsTeam               & 1{,}045{,}713 \\
Organization             & 1{,}043{,}512 \\
PostalAddress            & 1{,}030{,}249 \\
Country                  &   474{,}520 \\
AggregateRating          &   462{,}696 \\
SportsEvent              &   444{,}358 \\
Person                   &   410{,}829 \\
SiteNavigationElement    &   384{,}436 \\
BreadcrumbList           &   302{,}673 \\
Place                    &   580{,}546 \\
NewsArticle              &   199{,}809 \\
NewsMediaOrganization    &   184{,}805 \\
Article                  &   172{,}671 \\
BlogPosting              &   115{,}725 \\
WebSite                  &   126{,}518 \\
WebPage                  &   140{,}079 \\
LocalBusiness            &   214{,}643 \\
Thing                    &   104{,}639 \\
\bottomrule
\end{tabular}
\end{table}


% =============================================================================
\section{Full-Corpus Quality Analysis}
\label{sec:quality}
% =============================================================================

\subsection{From N-Quads to DuckDB}

Applying the in-memory entity-store approach to the full 14 GB merged file
caused an out-of-memory error.  The solution switches to a DuckDB~\citep{raasveldt2019duckdb} columnar
store:

\begin{enumerate}[leftmargin=1.5em]
  \item The 14 GB file is read \textbf{line by line} in a single Python process.
  \item Every \num{1000000} lines, four Python lists are flushed to a
        $\sim$100 MB Parquet file and immediately cleared.
  \item After all $\sim$70 Parquet files are written, DuckDB bulk-loads them
        using its own internal streaming; no large Python object is ever held
        in RAM.
\end{enumerate}

\begin{table}[H]
\centering
\caption{DuckDB entity store statistics.}
\label{tab:db}
\begin{tabular}{lr}
\toprule
\textbf{Metric} & \textbf{Value} \\
\midrule
Triples ingested        & 69{,}600{,}381 \\
Parse errors (skipped)  &     20{,}865 \; (0.03\%) \\
Entities                & 12{,}211{,}353 \\
DuckDB store size       & 18.61 GB \\
\bottomrule
\end{tabular}
\end{table}

The 20{,}865 parse errors (0.03\%) are malformed N-Quad lines falling into two
categories: \textbf{20{,}863} lines from \texttt{navbharattimes.indiatimes.com}
carry language-tagged literals with an unfilled template placeholder
(\texttt{"Hindi News"@\$lang}), where the site's JSON-LD renderer fails to
substitute \texttt{\$lang} with a valid BCP~47 language tag; and \textbf{2}
lines from \texttt{politifact.com} embed an HTML fragment inside a URI field
(e.g., \texttt{<https://schema.org/url</div>{}>}), making the predicate
unparseable.



\subsection{Namespace Correctness}

\begin{table}[H]
\centering
\caption{Schema.org namespace usage in the merged corpus.}
\label{tab:namespace}
\begin{tabular}{lrr}
\toprule
\textbf{URI class} & \textbf{Count} & \textbf{\%} \\
\midrule
Canonical (\texttt{http://schema.org/} or \texttt{https://schema.org/}) & 56{,}962{,}451 & 99.90\% \\
Non-canonical (\texttt{https://www.schema.org/})                         &     57{,}157   &  0.10\% \\
Truly malformed                                                           &           0    &  0.00\% \\
Other namespaces (\texttt{rdf:}, \texttt{owl:}, etc.)                    &    405{,}175   &    -- \\
\bottomrule
\end{tabular}
\end{table}

The top non-canonical example is
\texttt{https://www.schema.org/SiteNavigationElement} (17{,}372 occurrences);
schema.org does not use the \texttt{www.} subdomain.

\subsection{Type Validity}

Of all \texttt{rdf:type} assertions referencing Schema.org, \textbf{99.76\%}
use a class name that exists in the schema.org vocabulary (979 known classes).
Invalid types account for 28{,}102 declarations (0.24\%), predominantly
capitalisation errors or retired class names.
Table~\ref{tab:invalid_types} lists the top invalid type names.

\begin{table}[H]
\centering
\caption{Top invalid \texttt{rdf:type} names by declaration count.}
\label{tab:invalid_types}
\begin{tabular}{lr}
\toprule
\textbf{Invalid type} & \textbf{Declarations} \\
\midrule
imageObject          & 16{,}802 \\
PERSON               &  6{,}286 \\
person               &  2{,}500 \\
Website              &  1{,}015 \\
place                &    526 \\
Webpage              &    458 \\
liveBlogPosting      &    319 \\
webpage              &    161 \\
NewsPaper            &     15 \\
Expert               &      7 \\
\bottomrule
\end{tabular}
\end{table}


\subsection{Property Validity}

Property validity is assessed by checking whether each schema.org property is
actually defined for the declaring entity's type (using \texttt{domainIncludes}
with inheritance).

Table~\ref{tab:top_valid_props} lists the 20 most frequently used valid
Schema.org properties across the entire corpus (44{,}753{,}999 valid uses total).

\begin{table}[H]
\centering
\caption{Top 20 most-used valid Schema.org properties.}
\label{tab:top_valid_props}
\begin{tabular}{lr}
\toprule
\textbf{Property} & \textbf{Uses} \\
\midrule
name             & 6{,}578{,}093 \\
url              & 4{,}920{,}617 \\
height           & 1{,}644{,}399 \\
width            & 1{,}615{,}246 \\
position         & 1{,}512{,}667 \\
image            & 1{,}421{,}062 \\
address          & 1{,}283{,}189 \\
description      & 1{,}179{,}263 \\
contentUrl       & 1{,}104{,}876 \\
item             & 1{,}017{,}368 \\
addressLocality  & 1{,}005{,}532 \\
addressRegion    &   945{,}432   \\
logo             &   918{,}988   \\
addressCountry   &   807{,}141   \\
postalCode       &   761{,}799   \\
streetAddress    &   755{,}247   \\
telephone        &   728{,}709   \\
sameAs           &   672{,}483   \\
datePublished    &   616{,}006   \\
dateModified     &   613{,}182   \\
\bottomrule
\end{tabular}
\end{table}


\begin{table}[H]
\centering
\caption{Property validity breakdown.}
\label{tab:prop_validity}
\begin{tabular}{lrr}
\toprule
\textbf{Category} & \textbf{Count} & \textbf{\%} \\
\midrule
Valid property uses   & 44{,}753{,}999 & 99.18\% \\
Invalid property uses &    370{,}086   &  0.82\% \\
\midrule
\multicolumn{3}{l}{\small\itshape Breakdown of invalid uses by root cause:} \\
\midrule
Untyped entity (no \texttt{rdf:type})  & 115{,}149 & 31.1\% \\
Unknown entity type                    & 101{,}467 & 27.4\% \\
Unknown property (not in vocab)        &  93{,}779 & 25.3\% \\
Domain mismatch (wrong type)           &  59{,}691 & 16.1\% \\
\bottomrule
\end{tabular}
\end{table}

Property violations are classified into four mutually exclusive categories.
An \textbf{untyped entity} violation occurs when a triple's subject carries no
\texttt{rdf:type} declaration, making it impossible to verify whether the
property is applicable.
An \textbf{unknown type} violation arises when the declared entity type is not
recognised in the Schema.org vocabulary, so no valid property set can be
derived.
An \textbf{unknown property} violation is recorded when the property name
itself does not exist in the Schema.org vocabulary, regardless of the entity
type.
Finally, a \textbf{domain mismatch} violation occurs when both the entity type
and the property are individually valid Schema.org terms, but the property's
\texttt{domainIncludes} constraint (including inherited types) does not cover
the declared entity type.


The top-20 individual invalid property names broken down by violation reason are shown in Table~\ref{tab:invalid_props}.

\begin{table}[H]
\centering
\caption{Top 20 invalid property names by total misuse count, broken down by violation reason.}
\label{tab:invalid_props}
\begin{tabular}{lrrrrr}
\toprule
\textbf{Property} & \textbf{Total} & \textbf{Untyped} & \textbf{Unk.\ type} & \textbf{Unk.\ prop} & \textbf{Dom.\ mismatch} \\
\midrule
url                &  69{,}754 &  52{,}232 &  17{,}522 &       0 &          0 \\
name               &  69{,}619 &  59{,}167 &  10{,}452 &       0 &          0 \\
price              &  48{,}464 &      98   &       0   &       0 &  48{,}366 \\
query-input        &  24{,}424 &       0   &       0   &  24{,}424 &        0 \\
width              &  19{,}419 &       0   &  16{,}802 &       0 &   2{,}617 \\
description        &  19{,}219 &      12   &  19{,}207 &       0 &          0 \\
postID             &  17{,}323 &       0   &       0   &  17{,}323 &        0 \\
height             &  16{,}804 &       0   &  16{,}802 &       0 &          2 \\
Telephone          &  16{,}022 &       0   &       0   &  16{,}022 &        0 \\
caption            &  16{,}017 &       0   &  16{,}017 &       0 &          0 \\
partner            &  15{,}085 &       0   &       0   &  15{,}085 &        0 \\
@name              &   9{,}029 &   1{,}701 &       0   &   7{,}328 &        0 \\
primaryImageOfPage &   3{,}896 &       0   &       0   &       0 &   3{,}896 \\
SameAs             &   3{,}667 &       0   &       0   &   3{,}667 &        0 \\
jobTitle           &   1{,}820 &       0   &       7   &       0 &   1{,}813 \\
haspart            &   1{,}596 &       0   &       0   &   1{,}596 &        0 \\
publisher          &   1{,}120 &      12   &     323   &       0 &      785 \\
address            &   1{,}094 &     412   &     526   &       0 &      156 \\
isAccessibleForFree&     896   &       0   &     896   &       0 &          0 \\
potentialAction    &     896   &       0   &     896   &       0 &          0 \\
\bottomrule
\end{tabular}
\end{table}

The detailed domain-mismatch breakdown (entity type $\to$ property pairs) is shown in
Table~\ref{tab:domain_mismatch}.

\begin{table}[H]
\centering
\caption{Top 10 Domain-mismatch violations by entity type and property.}
\label{tab:domain_mismatch}
\begin{tabular}{llr}
\toprule
\textbf{Entity type} & \textbf{Property} & \textbf{Count} \\
\midrule
MedicalProcedure   & price               & 48{,}366 \\
Person             & width               &  2{,}615 \\
Article            & primaryImageOfPage  &  1{,}948 \\
NewsArticle        & primaryImageOfPage  &  1{,}948 \\
Organization       & jobTitle            &  1{,}813 \\
Organization       & publisher           &    769 \\
PerformanceRole    & actor               &    227 \\
TouristAttraction  & author              &    210 \\
TouristAttraction  & reviewBody          &    170 \\
Event              & address             &    156 \\
\bottomrule
\end{tabular}
\end{table}


\subsection{Range Violations}

Range violations occur when a property's value does not conform to the
expected type specified by its \texttt{rangeIncludes} constraint in the
Schema.org vocabulary.
For example, a property that expects a \texttt{URL} or \texttt{Text} literal
may instead receive a nested entity, or a property expecting an
\texttt{Integer} may receive a free-form string.
Unlike domain violations, which concern the \emph{subject} of a triple, range
violations concern the \emph{object}: the value provided is incompatible with
what the vocabulary prescribes for that property.


\begin{table}[H]
\centering
\caption{Top range-violating properties (merged corpus).}
\label{tab:range}
\begin{tabular}{lrl}
\toprule
\textbf{Property} & \textbf{Violations} & \textbf{Example values} \\
\midrule
\texttt{height}                & 1{,}640{,}370 & \texttt{"60"}, \texttt{"675"} (not QuantitativeValue) \\
\texttt{width}                 & 1{,}613{,}832 & \texttt{"600"}, \texttt{"300"} (not QuantitativeValue) \\
\texttt{item}                  &   611{,}303   & literal URL strings instead of URIs \\
\texttt{endDate}               &   395{,}515   & malformed dates / empty strings \\
\texttt{startDate}             &   382{,}002   & \texttt{"Sep 14, 2024"} (not ISO 8601) \\
\texttt{eventStatus}           &   188{,}068   & \texttt{"//schema.org/EventScheduled"} (malformed URI) \\
\texttt{eventAttendanceMode}   &   158{,}336   & \texttt{"MixedEventAttendanceMode"} (not a URI) \\
\texttt{cssSelector}           &   156{,}188   & CSS selector strings (internal use) \\
\texttt{associatedMedia}       &    95{,}569   & literal URL strings \\
\texttt{contentLocation}       &    94{,}625   & plain text location names (e.g.\ ``крым'') \\
\texttt{datePublished}         &    83{,}558   & non-ISO date strings \\
\texttt{mainEntity}            &    69{,}442   & Wikidata URI used as literal string \\
\texttt{xpath}                 &    65{,}520   & literal XPath selector strings \\
\texttt{author}                &    63{,}038   & plain string instead of Person/Organization \\
\texttt{dateModified}          &    60{,}531   & non-ISO date strings \\
\midrule
\multicolumn{3}{l}{\small Total checked: 55{,}056{,}887 assertions \quad Total violations: 6{,}054{,}576 (11.00\%)} \\
\bottomrule
\end{tabular}
\end{table}

The two largest violating properties (\texttt{height}, \texttt{width}) account
for 54\% of all violations: publishers routinely provide raw integers
where Schema.org expects a \texttt{QuantitativeValue} or \texttt{Distance} node.

\subsection{Distinct Schema.org Types per Page}

Across the 596{,}604 pages in the corpus, the mean number of distinct Schema.org
types declared is \textbf{7.01} (median~7.0, max~29).
Table~\ref{tab:types_per_page_dist} shows the distribution and
Table~\ref{tab:types_per_page_domain} highlights representative domains.

% maybe we can have this as a bar plot (maybe stacked bar plot that specifies the domains existing in each bar)

\begin{table}[H]
\centering
\caption{Distribution of distinct Schema.org type count per page.}
\label{tab:types_per_page_dist}
\begin{tabular}{rrr}
\toprule
\textbf{Types/page} & \textbf{Pages} & \textbf{\%} \\
\midrule
2  &  29{,}296 &  4.9\% \\
3  & 106{,}194 & 17.8\% \\
4  &  62{,}479 & 10.5\% \\
5  &  41{,}462 &  6.9\% \\
7  &  77{,}075 & 12.9\% \\
8  &  17{,}545 &  2.9\% \\
9  &  36{,}538 &  6.1\% \\
10 & 134{,}576 & 22.6\% \\
11 &  24{,}943 &  4.2\% \\
12+ & (remainder) & \\
\midrule
\textbf{Mean} & \multicolumn{2}{l}{7.01 \quad Median 7.0 \quad Max 29} \\
\bottomrule
\end{tabular}
\end{table}

% \begin{table}[H]
% \centering
% \caption{Selected domains: distinct types per page.}
% \label{tab:types_per_page_domain}
% \begin{tabular}{lrrr}
% \toprule
% \textbf{Domain} & \textbf{Pages} & \textbf{Avg types} & \textbf{Median} \\
% \midrule
% demagog.org.pl       &  1{,}335 & 14.87 & 15 \\
% freepressjournal.in  & 15{,}986 & 14.00 & 14 \\
% maariv.co.il         &  2{,}083 & 12.71 & 13 \\
% cbsnews.com          & 12{,}334 & 11.72 & 12 \\
% rbc.ua               & 30{,}542 & 10.00 & 10 \\
% ria.ru               &101{,}258 &  9.84 & 10 \\
% en.wikipedia.org     & 66{,}735 &  3.00 &  3 \\
% dw.com               &206{,}665 &  2.74 &  3 \\
% laverdad.es          & 15{,}518 &  2.00 &  2 \\
% \bottomrule
% \end{tabular}
% \end{table}

\subsection{Annotation Richness: Average Properties per Entity (AvP)}

The \textit{Average Properties per Entity} (AvP) metric quantifies the descriptive depth of structured annotations by measuring the mean number of distinct properties populated per entity instance. Formally, given a set of extracted entities $E$, AvP is defined as:

\begin{equation}
    \text{AvP} = \frac{1}{|E|} \sum_{e \in E} |P(e)|
\end{equation}

\noindent where $P(e)$ denotes the set of properties populated for entity $e$. A higher AvP indicates that publishers provide more detailed, semantically rich markup, whereas a low AvP suggests sparse annotation where only mandatory or minimal properties are used. In the context of \texttt{schema.org} types deployed across web pages, AvP serves as a proxy for annotation quality: an entity annotated with only a \texttt{url} and \texttt{name} contributes far less structured knowledge than one that additionally specifies \texttt{author}, \texttt{datePublished}, \texttt{claimReviewed}, and \texttt{reviewRating}. Comparing AvP across domains or entity types therefore reveals not only how broadly structured data is adopted, but also how thoroughly it is used.


% make it a bar plot
\begin{table}[H]
\centering
\caption{Corpus-wide AvP distribution.}
\label{tab:avp_dist}
\begin{tabular}{rrr}
\toprule
\textbf{Properties} & \textbf{Entities} & \textbf{\%} \\
\midrule
0  &   460{,}708 &  3.8\% \\
1  & 2{,}605{,}661 & 21.3\% \\
2  & 2{,}601{,}181 & 21.3\% \\
3  & 1{,}781{,}355 & 14.6\% \\
4  & 1{,}862{,}873 & 15.3\% \\
5  &   983{,}368 &  8.1\% \\
6  &   283{,}701 &  2.3\% \\
7  &   150{,}826 &  1.2\% \\
8  &   297{,}050 &  2.4\% \\
9  &   303{,}755 &  2.5\% \\
10 &   262{,}424 &  2.1\% \\
11 &   311{,}331 &  2.5\% \\
12--19 & 209{,}832 & 1.7\% \\
20+ &    96{,}348 &  0.8\% \\
\midrule
\textbf{Mean} & \multicolumn{2}{r}{3.70 \quad Median 3.0 \quad SD 3.60} \\
\bottomrule
\end{tabular}
\end{table}

\begin{table}[H]
\centering
\caption{AvP breakdown by Schema.org type for example types (full results for all 47 analysed types in Appendix).}
\label{tab:avp_bytype}
\begin{tabular}{lrrrr}
\toprule
\textbf{Type} & \textbf{Entities} & \textbf{Avail.\ props} & \textbf{AvP mean} & \textbf{Median} \\
\midrule
Article       &   172{,}671 & 136 & 19.57 & 25.0 \\
NewsArticle   &   199{,}809 & 141 & 12.61 & 14.0 \\
ClaimReview   &     8{,}729 & 139 &  6.60 &  6.0 \\
Event         &    79{,}947 &  56 &  6.58 &  5.0 \\
Organization  & 1{,}009{,}766 &  89 &  5.92 &  4.0 \\
Rating        &    11{,}486 &  19 &  2.88 &  3.0 \\
Product       &    45{,}456 &  72 &  2.81 &  2.0 \\
Person        &   410{,}829 &  81 &  2.79 &  2.0 \\
Claim         &     4{,}007 & 132 &  2.62 &  3.0 \\
Place         &   580{,}546 &  60 &  2.02 &  2.0 \\
\bottomrule
\end{tabular}
\end{table}

\subsection{Property Coverage per Type}

For each type, \emph{property coverage} is the average fraction of
schema.org-defined properties used per entity (entities using no
defined property contribute zero to the average). Across all 137
types in the corpus, the macro-averaged coverage is $5.75\%$
($\sigma = 4.64\%$), reflecting that most entities instantiate only
a small subset of their type's defined vocabulary.

\begin{table}[H]
\centering
\caption{Average per-entity property coverage for example types (full results for all 47 analysed types in Appendix).}
\label{tab:prop_cov}
\begin{tabular}{lrrr}
\toprule
\textbf{Type} & \textbf{Entities} & \textbf{Defined} & \textbf{Avg Coverage} \\
\midrule
Rating        &     11{,}516 &  19 & 15.09\% \\
Article       &    172{,}671 & 136 & 14.38\% \\
Event         &     80{,}561 &  56 & 11.63\% \\
NewsArticle   &    199{,}809 & 141 &  8.79\% \\
Organization  & 1{,}010{,}110 &  89 &  6.64\% \\
ClaimReview   &      8{,}729 & 139 &  4.74\% \\
Product       &     45{,}457 &  72 &  3.88\% \\
Person        &    410{,}961 &  81 &  3.44\% \\
Place         &    580{,}598 &  60 &  3.36\% \\
Claim         &      4{,}007 & 132 &  1.98\% \\
\midrule
\multicolumn{3}{l}{\textit{Macro avg (137 types)}} & 5.75\% \\
\multicolumn{3}{l}{\textit{Std dev}} & 4.64\% \\
\bottomrule
\end{tabular}
\end{table}



\subsection{Characteristic Property Sets}

A \emph{characteristic set} is the exact set of schema.org properties
used by an entity. Tables below show the three most frequent sets for
each of the ten example types; full results for all 47 analysed
types are provided in the Appendix.

\medskip\noindent\textbf{ClaimReview (24 unique sets).}

\begin{center}
\begin{tabular}{rlL{8cm}}
\toprule
\# & Count (\%) & Property set \\
\midrule
1 & 4{,}107 (47.1\%) & \{author, claimReviewed, datePublished, itemReviewed, reviewRating, url\} \\
2 & 1{,}719 (19.7\%) & \{author, claimReviewed, datePublished, description, itemReviewed, reviewRating, url\} \\
3 &   775   (8.9\%)  & \{author, claimReviewed, datePublished, reviewRating, url\} \\
\bottomrule
\end{tabular}
\end{center}

\medskip\noindent\textbf{Claim (13 unique sets).}

\begin{center}
\begin{tabular}{rlL{8cm}}
\toprule
\# & Count (\%) & Property set \\
\midrule
1 & 2{,}533 (63.2\%) & \{appearance, author, datePublished\} \\
2 &   520  (13.0\%)  & \{datePublished\} \\
3 &   410  (10.2\%)  & \{author, datePublished, url\} \\
\bottomrule
\end{tabular}
\end{center}

\medskip\noindent\textbf{NewsArticle (224 unique sets).}

\begin{center}
\begin{tabular}{rlL{8cm}}
\toprule
\# & Count (\%) & Property set \\
\midrule
1 & 41{,}330 (20.7\%) & \{author, dateModified, datePublished, headline, image, mainEntityOfPage, publisher\} \\
2 & 18{,}582  (9.3\%) & \{articleBody, articleSection, author, dateModified, datePublished, description, hasPart, headline, image, inLanguage, isAccessibleForFree, isPartOf, keywords, mainEntityOfPage, publisher, url\} \\
3 & 17{,}602  (8.8\%) & \{articleBody, articleSection, author, dateCreated, dateModified, datePublished, description, headline, image, keywords, locationCreated, mainEntityOfPage, publisher, publishingPrinciples\} \\
\bottomrule
\end{tabular}
\end{center}

\medskip\noindent\textbf{Article (133 unique sets).}

\begin{center}
\begin{tabular}{rlL{8cm}}
\toprule
\# & Count (\%) & Property set \\
\midrule
1 & 49{,}185 (28.5\%) & \{author, dateModified, datePublished, headline, image, mainEntity, name, publisher, sameAs, url\} \\
2 & 17{,}949 (10.4\%) & \{author, dateModified, datePublished, headline, mainEntity, name, publisher, sameAs, url\} \\
3 & 16{,}179  (9.4\%) & \{alternativeHeadline, associatedMedia, author, copyrightHolder, copyrightYear, creator, dateCreated, dateModified, datePublished, description, discussionUrl, genre, headline, image, inLanguage, keywords, license, mainEntityOfPage, name, publisher, publishingPrinciples, sourceOrganization, speakable, thumbnailUrl, url\} \\
\bottomrule
\end{tabular}
\end{center}

\medskip\noindent\textbf{Organization (63 unique sets).}

\begin{center}
\begin{tabular}{rlL{8cm}}
\toprule
\# & Count (\%) & Property set \\
\midrule
1 & 295{,}728 (29.3\%) & \{address, aggregateRating, award, email, image, legalName, logo, name, sameAs, telephone, url\} \\
2 & 215{,}329 (21.3\%) & \{name\} \\
3 & 166{,}687 (16.5\%) & \{logo, name\} \\
\bottomrule
\end{tabular}
\end{center}

\medskip\noindent\textbf{Person (170 unique sets).}

\begin{center}
\begin{tabular}{rlL{8cm}}
\toprule
\# & Count (\%) & Property set \\
\midrule
1 & 128{,}766 (31.3\%) & \{name, url\} \\
2 & 124{,}946 (30.4\%) & \{name\} \\
3 &  55{,}459 (13.5\%) & \{image, name, url\} \\
\bottomrule
\end{tabular}
\end{center}

\medskip\noindent\textbf{Event (52 unique sets).}

\begin{center}
\begin{tabular}{rlL{8cm}}
\toprule
\# & Count (\%) & Property set \\
\midrule
1 & 47{,}077 (58.4\%) & \{endDate, location, name, startDate, url\} \\
2 & 29{,}364 (36.4\%) & \{description, endDate, eventStatus, image, location, name, performer, startDate, url\} \\
3 &    614   (0.8\%)  & \{(no SO properties)\} \\
\bottomrule
\end{tabular}
\end{center}

\medskip\noindent\textbf{Place (18 unique sets).}

\begin{center}
\begin{tabular}{rlL{8cm}}
\toprule
\# & Count (\%) & Property set \\
\midrule
1 & 519{,}513 (89.5\%) & \{address, name\} \\
2 &  28{,}241  (4.9\%) & \{name\} \\
3 &  14{,}241  (2.5\%) & \{address, geo, hasMap, name, url\} \\
\bottomrule
\end{tabular}
\end{center}

\medskip\noindent\textbf{Product (31 unique sets).}

\begin{center}
\begin{tabular}{rlL{8cm}}
\toprule
\# & Count (\%) & Property set \\
\midrule
1 & 38{,}792 (85.3\%) & \{name, productID\} \\
2 &  3{,}048  (6.7\%) & \{brand, category, description, image, name, offers, sku, url\} \\
3 &    849   (1.9\%)  & \{brand, category, description, image, name, offers, releaseDate, sku, url\} \\
\bottomrule
\end{tabular}
\end{center}

\medskip\noindent\textbf{Rating (9 unique sets).}

\begin{center}
\begin{tabular}{rlL{8cm}}
\toprule
\# & Count (\%) & Property set \\
\midrule
1 & 4{,}104 (35.6\%) & \{alternateName, bestRating, ratingValue, worstRating\} \\
2 & 3{,}603 (31.3\%) & \{alternateName\} \\
3 & 2{,}593 (22.5\%) & \{bestRating, ratingValue, worstRating\} \\
\bottomrule
\end{tabular}
\end{center}


\subsection{Property Fill Rate by Domain}

Fill rate per entity = distinct properties used / properties defined for the
entity's type.  Corpus-wide mean fill rate: \textbf{7.4\%}
(11{,}823{,}497 typed entities; 45{,}124{,}082 distinct entity--predicate pairs).
Domain-level aggregation is at the registered-domain level (e.g.\ subdomains of
\texttt{cnn.com} are merged).

\begin{table}[H]
\centering
\small
\setlength{\tabcolsep}{4pt}
\caption{Property fill rate for selected fact-checking organisations, news outlets,
  and knowledge bases, aggregated at the registered-domain level,
  sorted by fill rate within each group (descending).
  \textit{Entities}: typed entities from this domain;
  \textit{Avail.}: avg. number of Schema.org properties applicable to the entity's declared type;
  \textit{Avg Used}: mean properties used per entity;
  \textit{Fill\%}: avg used / avail.}
\label{tab:fillrate}
\begin{tabular}{lrrrr}
\toprule
\textbf{Domain} & \textbf{Entities} & \textbf{Avail.} & \textbf{Avg Used} & \textbf{Fill\%} \\
\midrule
\multicolumn{5}{l}{\textit{Fact-checking organisations}} \\
\midrule
politifact.com        &   3{,}935  &  91 & 4.10 &  8.0\% \\
afp.com               &   2{,}373  &  93 & 3.61 &  6.6\% \\
vishvasnews.com       &  16{,}987  &  85 & 3.57 &  5.8\% \\
facta.news            &   8{,}631  &  70 & 3.12 &  5.8\% \\
demagog.org.pl        &  23{,}075  &  75 & 3.05 &  5.8\% \\
maldita.es            &   6{,}094  & 107 & 3.13 &  4.8\% \\
factcheck.org         &  17{,}539  &  96 & 2.70 &  4.8\% \\
leadstories.com       &  17{,}200  & 108 & 3.40 &  4.5\% \\
fullfact.org          &   9{,}442  &  97 & 2.86 &  3.2\% \\
\midrule
\multicolumn{5}{l}{\textit{News outlets (global)}} \\
\midrule
cbsnews.com           &  249{,}990  &  99 & 3.80 &  5.6\% \\
cnn.com               &  281{,}214  & 113 & 5.02 &  4.3\% \\
aljazeera.com         &    2{,}971  & 117 & 2.97 &  2.9\% \\
dw.com                &  208{,}338  & 121 & 2.74 &  2.3\% \\
\midrule
\multicolumn{5}{l}{\textit{News outlets (national/regional)}} \\
\midrule
lefigaro.fr           &   339{,}409   &  45 & 3.20 & 10.7\% \\
ria.ru                & 3{,}537{,}108 &  78 & 5.08 & 10.0\% \\
elpais.com            & 1{,}148{,}493 &  53 & 2.85 &  9.6\% \\
rbc.ua                &   408{,}030   &  62 & 4.20 &  9.3\% \\
elcomercio.pe         &    32{,}489   & 114 & 6.94 &  6.6\% \\
lequipe.fr            &    11{,}259   &  83 & 3.63 &  6.6\% \\
lavanguardia.com      &    18{,}085   & 119 & 7.24 &  6.1\% \\
clarin.com            &     5{,}193   &  86 & 3.70 &  5.9\% \\
indiatimes.com        &   499{,}882   & 104 & 3.39 &  4.5\% \\
freepressjournal.in   &   659{,}227   & 109 & 2.88 &  4.3\% \\
pravda.ru             &    57{,}486   &  36 & 1.51 &  4.2\% \\
infobae.com           &    70{,}679   & 128 & 3.19 &  2.6\% \\
\midrule
\multicolumn{5}{l}{\textit{Knowledge base}} \\
\midrule
wikipedia.org         &   277{,}768   & 116 & 3.42 &  2.8\% \\
\bottomrule
\end{tabular}
\end{table}

\begin{table}[H]
\centering
\caption{Distribution of per-entity fill rates across the full typed-entity corpus
  (11{,}823{,}497 entities).}
\label{tab:fillrate_dist}
\begin{tabular}{lrr}
\toprule
\textbf{Fill rate range} & \textbf{Entities} & \textbf{\%} \\
\midrule
0--5\%    & 6{,}344{,}431 & 53.7\% \\
5--10\%   & 1{,}442{,}584 & 12.2\% \\
10--20\%  & 3{,}647{,}003 & 30.8\% \\
20--30\%  &   389{,}462   &  3.3\% \\
30--50\%  &        16     &  \textless{}0.1\% \\
50--100\% &         1     &  \textless{}0.1\% \\
\bottomrule
\end{tabular}
\end{table}

\subsection{Article and NewsArticle Pages}

News and article pages are primary candidate sources for
automated fact-checking: they carry the claims, contextual evidence, and
source references that fact-checking pipelines need to reason over.
To understand how much structured semantic information Schema.org annotations
already provide on such pages, and therefore what LLM-based enrichment would
still need to supply, we identify a representative set of article and
news-article pages from the corpus and analyse the richness and coverage of
their existing annotations.

Pages were identified using four complementary signals: single entity (Article and/or NewsArticle) per
page, subject URI matching the page URI, presence of
\texttt{headline} + \texttt{datePublished}/author/body, and URL path
heuristics.

\begin{table}[H]
\centering
\caption{Article page identification.}
\label{tab:article_id}
\begin{tabular}{lr}
\toprule
\textbf{Filter} & \textbf{Pages} \\
\midrule
Candidate pages (Article or NewsArticle declared) & 356{,}132 \\
Signal 1: single entity per page & 342{,}808 \\
Signal 2: subject URL matches page URL & 5{,}613 \\
Signal 3: has headline + date/author/body & 354{,}084 \\
Signal 4: URL path looks like article & 354{,}740 \\
\midrule
Strict intersection (all 4 signals) & 5{,}610 \\
\textbf{Lenient intersection (signals 1+3)} & \textbf{340{,}789} \\
\bottomrule
\end{tabular}
\end{table}

\begin{table}[H]
\centering
\caption{Article and NewsArticle statistics on 340{,}789 lenient article pages.}
\label{tab:articles}
\begin{tabular}{lrrrr}
\toprule
\textbf{Type} & \textbf{Entities} & \textbf{AvP} & \textbf{Avail.\ props} & \textbf{No-prop \%} \\
\midrule
Article     & 158{,}695 & 20.1 & 136 & 0.0\% \\
NewsArticle & 178{,}854 & 12.6 & 141 & 0.0\% \\
Overlap     &      81   & -- & -- & -- \\
\bottomrule
\end{tabular}
\end{table}

Beyond overall annotation depth, we examine four Schema.org properties with
direct semantic relevance to fact-checking: \texttt{articleBody} (full article
text), \texttt{about} (the
entity the article reports on), \texttt{citation} (external sources cited by the author, supporting
source verification), and \texttt{mentions} (named entities referenced in the
text, which if populated would provide structured named-entity annotations
without requiring additional extraction from raw text).
Table~\ref{tab:article_props} reports their actual fill rates.



\begin{table}[H]
\centering
\caption{Content property fill rates on article pages.}
\label{tab:article_props}
\begin{tabular}{lrrrl}
\toprule
\textbf{Property} & \textbf{Article} & \textbf{NewsArticle} & \textbf{Combined} & \textbf{Note} \\
\midrule
\texttt{articleBody} & 32.2\%  & 61.6\% & 47.8\% & full text \\
\texttt{about}       & 39.6\%  &  1.3\% & 19.3\% & topic entity \\
\texttt{citation}    & 23.3\%  &  0.3\% & 11.1\% & references \\
\texttt{mentions}    & $<$0.1\% & $<$0.1\% & $<$0.1\% & named entities \\
\bottomrule
\end{tabular}
\end{table}

\noindent\textit{NewsArticle detail (178{,}854 entities):}
\texttt{articleBody} 110{,}241 (61.6\%);
\texttt{about} 2{,}286 (1.3\%);
\texttt{citation} 448 (0.3\%);
\texttt{mentions} 27 ($<$0.1\%).

\paragraph{Structure of \texttt{about} targets.}
Of the 65{,}146 Article and NewsArticle entities with an \texttt{about} link
(19.3\%), the linked targets comprise 63{,}411 distinct entities across 9 types.
Among those entities:
\begin{itemize}[leftmargin=1.5em]
  \item \texttt{Thing}: 62{,}955 (99.3\%): bare entities with no further
        structured properties.
  \item \texttt{SportsTeam}: 190 (0.3\%)
  \item \texttt{SportsEvent}: 180 (0.3\%)
  \item \texttt{Event}: 58 (0.1\%)
  \item \texttt{GovernmentOrganization}: 12; \texttt{Person}: 8;
        \texttt{EducationalOrganization}: 4; \texttt{Country}: 2;
        \texttt{Organization}: 2
\end{itemize}
Sample \texttt{name} values from \texttt{Thing} targets reveal that these are
not structured entities but free-text topic strings, typically comma-separated
keyword lists, providing
no usable structure.


\paragraph{Characteristic property sets.}
To understand \emph{what} is annotated on these pages and how richly, we
examine the characteristic property sets: distinct combinations of Schema.org
properties assigned to entities.
Table~\ref{tab:charsets_all} shows the top-3 most common sets for all 20
entity types on article pages: \texttt{Article}, \texttt{NewsArticle},
and the 18 co-occurring types catalogued in Table~\ref{tab:cooccur}.

\begin{table}[H]
\centering
\caption{Co-occurring entity types on Article/NewsArticle pages
(base: 340{,}789 lenient article pages).}
\label{tab:cooccur}
\begin{tabular}{lrr}
\toprule
\textbf{Type} & \textbf{Pages} & \textbf{\% of article pages} \\
\midrule
ImageObject              & 340{,}654 & 100.0\% \\
Organization             & 297{,}085 &  87.2\% \\
WebPage                  & 204{,}916 &  60.1\% \\
NewsArticle              & 182{,}175 &  53.5\% \\
PostalAddress            & 156{,}314 &  45.9\% \\
NewsMediaOrganization    & 156{,}825 &  46.0\% \\
Article                  & 158{,}695 &  46.6\% \\
ListItem                 & 168{,}463 &  49.4\% \\
BreadcrumbList           & 168{,}208 &  49.4\% \\
Person                   & 137{,}195 &  40.3\% \\
WebSite                  & 112{,}988 &  33.2\% \\
AggregateRating          &  96{,}115 &  28.2\% \\
Country                  &  95{,}034 &  27.9\% \\
Thing                    &  78{,}889 &  23.1\% \\
SpeakableSpecification   &  76{,}483 &  22.4\% \\
CreativeWork             &  50{,}186 &  14.7\% \\
AdministrativeArea       &  48{,}074 &  14.1\% \\
ContactPoint             &  44{,}991 &  13.2\% \\
WebPageElement           &  34{,}852 &  10.2\% \\
Product                  &  32{,}403 &   9.5\% \\
Place                    &  27{,}458 &   8.1\% \\
VideoObject              &  22{,}211 &   6.5\% \\
CollectionPage           &  12{,}836 &   3.8\% \\
\bottomrule
\end{tabular}
\end{table}

\setlength{\tabcolsep}{4pt}
\begin{longtable}{rrrp{9cm}}
\caption{Top-3 characteristic property sets for all entity types on article
  pages (340{,}789 lenient article pages).
  \% = fraction of entities of that type within the article-page corpus.}
\label{tab:charsets_all}\\
\toprule
\multicolumn{1}{c}{\textbf{\#}} & \textbf{Entities} & \textbf{\%} & \textbf{Properties} \\
\midrule
\endfirsthead
\multicolumn{4}{c}{\tablename~\ref{tab:charsets_all} (continued)}\\
\toprule
\multicolumn{1}{c}{\textbf{\#}} & \textbf{Entities} & \textbf{\%} & \textbf{Properties} \\
\midrule
\endhead
\midrule
\multicolumn{4}{r}{\small\itshape continued\ldots}\\
\endfoot
\bottomrule
\endlastfoot
%
\multicolumn{4}{l}{\textit{Article} (158{,}695 entities)} \\
\midrule
1 & 45{,}131 & 28.4 & \scriptsize author, dateModified, datePublished, headline, image, mainEntity, name, publisher, sameAs, url \\[2pt]
2 & 17{,}393 & 11.0 & \scriptsize author, dateModified, datePublished, headline, mainEntity, name, publisher, sameAs, url \\[2pt]
3 & 15{,}905 & 10.0 & \scriptsize alternativeHeadline, associatedMedia, author, copyrightHolder, copyrightYear, creator, dateCreated, dateModified, datePublished, description, discussionUrl, genre, headline, image, inLanguage, keywords, license, mainEntityOfPage, name, publisher, publishingPrinciples, sourceOrganization, speakable, thumbnailUrl, url \\[4pt]
%
\multicolumn{4}{l}{\textit{NewsArticle} (178{,}854 entities)} \\
\midrule
1 & 40{,}755 & 22.8 & \scriptsize author, dateModified, datePublished, headline, image, mainEntityOfPage, publisher \\[2pt]
2 & 18{,}450 & 10.3 & \scriptsize articleBody, articleSection, author, dateModified, datePublished, description, hasPart, headline, image, inLanguage, isAccessibleForFree, isPartOf, keywords, mainEntityOfPage, publisher, url \\[2pt]
3 & 17{,}548 &  9.8 & \scriptsize articleBody, articleSection, author, dateCreated, dateModified, datePublished, description, headline, image, keywords, locationCreated, mainEntityOfPage, publisher, publishingPrinciples \\[4pt]
%
\multicolumn{4}{l}{\textit{ImageObject} (1{,}561{,}640 entities)} \\
\midrule
1 & 767{,}153 & 49.1 & \scriptsize contentUrl, height, url, width \\[2pt]
2 & 392{,}827 & 25.2 & \scriptsize height, url, width \\[2pt]
3 & 189{,}519 & 12.1 & \scriptsize contentUrl, height, representativeOfPage, url, width \\[4pt]
%
\multicolumn{4}{l}{\textit{Organization} (859{,}726 entities)} \\
\midrule
1 & 294{,}265 & 34.2 & \scriptsize address, aggregateRating, award, email, image, legalName, logo, name, sameAs, telephone, url \\[2pt]
2 & 162{,}374 & 18.9 & \scriptsize name \\[2pt]
3 & 158{,}699 & 18.5 & \scriptsize address, aggregateRating, award, email, legalName, logo, name, sameAs, telephone, url \\[4pt]
%
\multicolumn{4}{l}{\textit{ListItem} (568{,}698 entities)} \\
\midrule
1 & 284{,}856 & 50.1 & \scriptsize item, name, position, url \\[2pt]
2 & 147{,}397 & 25.9 & \scriptsize item, position \\[2pt]
3 & 130{,}009 & 22.9 & \scriptsize item, name, position \\[4pt]
%
\multicolumn{4}{l}{\textit{PostalAddress} (540{,}246 entities)} \\
\midrule
1 & 470{,}697 & 87.1 & \scriptsize addressCountry, addressLocality, addressRegion, postalCode, streetAddress \\[2pt]
2 &  32{,}460 &  6.0 & \scriptsize addressCountry, addressLocality, postalCode, streetAddress \\[2pt]
3 &  17{,}007 &  3.1 & \scriptsize \textit{(no properties)} \\[4pt]
%
\multicolumn{4}{l}{\textit{AggregateRating} (454{,}443 entities)} \\
\midrule
1 & 453{,}015 & 99.7 & \scriptsize bestRating, ratingCount, ratingValue, worstRating \\[2pt]
2 &   1{,}369 &  0.3 & \scriptsize bestRating, ratingCount, ratingValue \\[2pt]
3 &        45 & $<$0.1 & \scriptsize ratingCount, ratingValue \\[4pt]
%
\multicolumn{4}{l}{\textit{Country} (453{,}044 entities)} \\
\midrule
1 & 453{,}039 & 100.0 & \scriptsize name \\[2pt]
2 &         2 & $<$0.1 & \scriptsize description, name, url \\[2pt]
3 &         1 & $<$0.1 & \scriptsize item, name, position \\[4pt]
%
\multicolumn{4}{l}{\textit{Person} (210{,}433 entities)} \\
\midrule
1 &  81{,}050 & 38.5 & \scriptsize name, url \\[2pt]
2 &  75{,}654 & 36.0 & \scriptsize name \\[2pt]
3 &  16{,}444 &  7.8 & \scriptsize image, jobTitle, name, sameAs, url \\[4pt]
%
\multicolumn{4}{l}{\textit{WebPage} (204{,}916 entities)} \\
\midrule
1 &  88{,}240 & 43.1 & \scriptsize \textit{(no properties)} \\[2pt]
2 &  31{,}043 & 15.1 & \scriptsize name \\[2pt]
3 &  27{,}487 & 13.4 & \scriptsize name, speakable, url \\[4pt]
%
\multicolumn{4}{l}{\textit{BreadcrumbList} (170{,}358 entities)} \\
\midrule
1 & 169{,}694 & 99.6 & \scriptsize itemListElement \\[2pt]
2 &       664 &  0.4 & \scriptsize itemListElement, name \\[4pt]
%
\multicolumn{4}{l}{\textit{NewsMediaOrganization} (156{,}825 entities)} \\
\midrule
1 &  53{,}871 & 34.4 & \scriptsize logo, name \\[2pt]
2 &  37{,}564 & 24.0 & \scriptsize logo, name, url \\[2pt]
3 &  30{,}114 & 19.2 & \scriptsize foundingDate, logo, name, sameAs, url \\[4pt]
%
\multicolumn{4}{l}{\textit{WebSite} (102{,}011 entities)} \\
\midrule
1 &  94{,}955 & 93.1 & \scriptsize description, image, name, url \\[2pt]
2 &   3{,}843 &  3.8 & \scriptsize potentialAction, url \\[2pt]
3 &   3{,}179 &  3.1 & \scriptsize name, url \\[4pt]
%
\multicolumn{4}{l}{\textit{Thing} (79{,}127 entities)} \\
\midrule
1 &  75{,}242 & 95.1 & \scriptsize name \\[2pt]
2 &   3{,}756 &  4.7 & \scriptsize name, url \\[2pt]
3 &       124 &  0.2 & \scriptsize name, sameAs \\[4pt]
%
\multicolumn{4}{l}{\textit{SpeakableSpecification} (76{,}557 entities)} \\
\midrule
1 &  44{,}072 & 57.6 & \scriptsize cssSelector \\[2pt]
2 &  32{,}424 & 42.4 & \scriptsize xpath \\[2pt]
3 &        46 &  0.1 & \scriptsize cssSelector, xPath \\[4pt]
%
\multicolumn{4}{l}{\textit{AdministrativeArea} (57{,}499 entities)} \\
\midrule
1 &  57{,}493 & 100.0 & \scriptsize name \\[2pt]
2 &         6 & $<$0.1 & \scriptsize description, name, url \\[4pt]
%
\multicolumn{4}{l}{\textit{ContactPoint} (54{,}341 entities)} \\
\midrule
1 &  29{,}616 & 54.5 & \scriptsize areaServed, availableLanguage, contactType, email, telephone \\[2pt]
2 &  17{,}366 & 32.0 & \scriptsize contactType, telephone \\[2pt]
3 &   6{,}522 & 12.0 & \scriptsize contactType, email \\[4pt]
%
\multicolumn{4}{l}{\textit{CreativeWork} (51{,}331 entities)} \\
\midrule
1 &  31{,}539 & 61.4 & \scriptsize name, productID \\[2pt]
2 &  16{,}019 & 31.2 & \scriptsize name, url \\[2pt]
3 &   1{,}707 &  3.3 & \scriptsize author, datePublished, keywords \\[4pt]
%
\multicolumn{4}{l}{\textit{WebPageElement} (34{,}852 entities)} \\
\midrule
1 &  34{,}852 & 100.0 & \scriptsize cssSelector, isAccessibleForFree \\[4pt]
%
\multicolumn{4}{l}{\textit{Product} (32{,}627 entities)} \\
\midrule
1 &  31{,}539 & 96.7 & \scriptsize name, productID \\[2pt]
2 &       551 &  1.7 & \scriptsize Name, brand, description, image, mpn, productID, sku \\[2pt]
3 &       241 &  0.7 & \scriptsize name \\
\end{longtable}
\setlength{\tabcolsep}{6pt}

\noindent
Across all 20 types, annotation follows highly regular, publisher-driven
templates, and the dominant property sets are oriented towards publication
metadata rather than semantic content.
\texttt{Article} entities are the richest in the corpus (AvP 20.1): all
characteristic sets share a stable publication-metadata core
(\texttt{author}, \texttt{datePublished}, \texttt{headline}, \texttt{image},
\texttt{publisher}, \texttt{url}), and richer sets extend this with editorial
detail (\texttt{genre}, \texttt{keywords}, \texttt{license},
\texttt{description}) and structural references such as \texttt{articleBody}
and \texttt{citation}, though even these remain metadata-adjacent rather than
content-grounding.
\texttt{NewsArticle} entities are more compact (AvP 12.6): the dominant set
(22.8\%) encodes just 7 metadata properties, and richer sets add paywall
signals (\texttt{isAccessibleForFree}, \texttt{isPartOf}, \texttt{hasPart})
and location metadata (\texttt{locationCreated}).
The co-occurring types reinforce the same template-driven pattern:
\texttt{BreadcrumbList} and \texttt{AggregateRating} are effectively
single-set types (99.6\% and 99.7\% respectively); geographical references
(\texttt{Country}, \texttt{AdministrativeArea}) converge at over 99\% on
the single-property set \{\texttt{name}\}; \texttt{ImageObject} (present on
virtually every article page) is annotated almost exclusively with technical
dimensions (\texttt{contentUrl}, \texttt{height}, \texttt{url},
\texttt{width}), with no caption, rights, or authorship metadata;
\texttt{WebPage} is the sparsest type, with 43.1\% of instances carrying no
Schema.org properties at all; and \texttt{Person} entities are thin, with
74.5\% annotated using only \{\texttt{name}\} or \{\texttt{name},
\texttt{url}\}.
Against this backdrop, the content properties examined in
Table~\ref{tab:article_props} are almost entirely absent: \texttt{mentions}
fills at effectively 0\%, \texttt{about} at 19.3\% with 99.3\% of its
targets being bare \texttt{Thing} nodes with no further structured
properties, and \texttt{citation} at 11.1\%.
Structured data on article pages therefore describes publication
infrastructure rather than semantic content, motivating the use of LLMs to
extract the richer, content-grounded annotations that Schema.org publishers
currently omit.

% =============================================================================
\section{ClaimReview Analysis}
\label{sec:claimreview}
% =============================================================================

This section reports on the full WDC 2024-12 ClaimReview class-specific
subset (not restricted to the target domain list), enabling a broad assessment
of how fact-checking organisations use Schema.org~\citep{schemaorg_claimreview}.
%Related work on knowledge graphs derived from ClaimReview data includes \citep{tchechmedjiev2019claimskg}.

\subsection{Dataset}

\begin{table}[H]
\centering
\caption{ClaimReview dataset overview.}
\label{tab:cr_dataset}
\begin{tabular}{lr}
\toprule
\textbf{Metric} & \textbf{Value} \\
\midrule
Source file (\texttt{part\_0.gz}) & 57.6 MB compressed \\
N-Quad lines parsed & 3{,}919{,}703 \\
Parse errors & 2 \\
All entities parsed & 787{,}354 \\
\textbf{ClaimReview entities} & \textbf{60{,}313} \\
Registered domains & 318 \\
\midrule
\multicolumn{2}{l}{\small\itshape Schema.org vocabulary for ClaimReview:} \\
Properties defined for ClaimReview & 139 \\
Properties actually used & 41 (29.5\%) \\
Properties never used & 98 (70.5\%) \\
\bottomrule
\end{tabular}
\end{table}

\subsection{Syntactic Validity}

\begin{table}[H]
\centering
\caption{ClaimReview syntactic quality metrics.}
\label{tab:cr_quality}
\begin{tabular}{lrr}
\toprule
\textbf{Metric} & \textbf{Value} & \textbf{\%} \\
\midrule
Namespace correctness  & 316{,}261 / 316{,}261 & 100.00\% \\
Type validity          &  60{,}539 /  60{,}540 & 100.00\% \\
Property validity      & 255{,}057 / 255{,}721 & 99.74\% \\
Range violations       &   5{,}858 / 258{,}093 &  2.27\% \\
\bottomrule
\end{tabular}
\end{table}

ClaimReview is significantly cleaner than the full-corpus average on all
four metrics.  The single type invalidity is \texttt{ExplainerNewsArticle}
(a non-standard extension type used by one publisher).  Range violations are
dominated by non-ISO date formats in \texttt{datePublished}
(4{,}990 cases; e.g., \texttt{"10/08/2023"}, \texttt{"3 Nov. 2018"}) and
strings in \texttt{potentialAction} where an \texttt{Action} node is expected
(488 cases).

\subsection{Property Coverage}

\begin{table}[H]
\centering
\caption{Top 20 used ClaimReview properties, sorted by entity density
  (density = fraction of ClaimReview entities that populate the property).}

\label{tab:cr_coverage}
\begin{tabular}{lrr}
\toprule
\textbf{Property} & \textbf{Entities using} & \textbf{Density \%} \\
\midrule
url                  & 36{,}466 & 60.46\% \\
author               & 36{,}274 & 60.14\% \\
claimReviewed        & 36{,}020 & 59.72\% \\
reviewRating         & 35{,}766 & 59.30\% \\
itemReviewed         & 33{,}730 & 55.92\% \\
datePublished        & 32{,}992 & 54.70\% \\
dateModified         &  6{,}813 & 11.30\% \\
image                &  6{,}306 & 10.46\% \\
publisher            &  5{,}649 &  9.37\% \\
mainEntityOfPage     &  5{,}278 &  8.75\% \\
name                 &  5{,}049 &  8.37\% \\
description          &  4{,}949 &  8.21\% \\
keywords             &  1{,}366 &  2.26\% \\
inLanguage           &  1{,}345 &  2.23\% \\
headline             &  1{,}291 &  2.14\% \\
reviewBody           &  1{,}239 &  2.05\% \\
sdPublisher          &    879   &  1.46\% \\
isPartOf             &    735   &  1.22\% \\
potentialAction      &    718   &  1.19\% \\
dateCreated          &    385   &  0.64\% \\
\bottomrule
\end{tabular}
\end{table}



Approximately 40\% of ClaimReview entities lack three core semantic
properties: \texttt{claimReviewed}, \texttt{reviewRating}, \texttt{itemReviewed}.
The \texttt{vpravda.ru} cluster (Section~\ref{sec:cr:domains}) accounts for
most of this gap.

\subsection{Annotation Richness (AvP)}

\begin{table}[H]
\centering
\caption{ClaimReview AvP distribution.}
\label{tab:cr_avp}
\begin{tabular}{rrr}
\toprule
\textbf{Properties} & \textbf{Entities} & \textbf{\%} \\
\midrule
0  & 22{,}896 & 38.0\% \\
1  &    165   &  0.3\% \\
2  &     14   &  0.0\% \\
3  &    161   &  0.3\% \\
4  &    976   &  1.6\% \\
5  &  5{,}587 &  9.3\% \\
6  & 17{,}220 & 28.5\% \\
7  &  5{,}474 &  9.1\% \\
8  &  1{,}524 &  2.5\% \\
9  &    492   &  0.8\% \\
10 &  4{,}359 &  7.2\% \\
11 &    475   &  0.8\% \\
15 &    633   &  1.0\% \\
20+ &   269   &  0.4\% \\
\midrule
\textbf{Mean 4.24} & \multicolumn{2}{l}{Median 6.0 \quad SD 3.79 \quad Max 24} \\
\bottomrule
\end{tabular}
\end{table}

\subsection{Nesting Depth}

\begin{table}[H]
\centering
\caption{ClaimReview nesting depth distribution.}
\label{tab:cr_depth}
\begin{tabular}{lrr}
\toprule
\textbf{Depth} & \textbf{Entities} & \textbf{\%} \\
\midrule
0 & 23{,}552 & 39.0\% \\
1 &  1{,}843 &  3.1\% \\
2 & 24{,}635 & 40.8\% \\
3 &  1{,}454 &  2.4\% \\
4 &  7{,}583 & 12.6\% \\
5 &    845   &  1.4\% \\
6 &    301   &  0.5\% \\
7 &      3   &  0.0\% \\
8 &     97   &  0.2\% \\
\midrule
\textbf{Mean} & \multicolumn{2}{l}{1.54 ($\sigma = 1.47$) \quad Median 2.0 \quad Max 8} \\
\midrule
Flat (depth = 0) & 23{,}552 & 39.0\% \\
Nested (depth $\geq 1$) & 36{,}761 & 61.0\% \\
Deep (depth $\geq 3$) & 10{,}283 & 17.0\% \\
\bottomrule
\end{tabular}
\end{table}

\begin{table}[H]
\centering
\caption{Cross-reference density for ClaimReview entities.}
\label{tab:cr_xref}
\begin{tabular}{lr}
\toprule
\textbf{Metric} & \textbf{Value} \\
\midrule
Total cross-refs             & 143{,}231 \\
Avg refs per entity          & 2.37 \\
Blank-node refs              & 115{,}168 (80.4\%) \\
URI refs                     &  28{,}063 (19.6\%) \\
\midrule
\multicolumn{2}{l}{\small\itshape Types of directly referenced entities:} \\
Organization  & 37{,}016 \\
Rating        & 35{,}115 \\
WebPage       & 21{,}994 \\
Claim         & 19{,}824 \\
CreativeWork  & 13{,}961 \\
ImageObject   &  5{,}704 \\
Person        &  5{,}566 \\
(untyped)     &  1{,}706 \\
LocalBusiness &  1{,}256 \\
ItemPage      &    946   \\
\bottomrule
\end{tabular}
\end{table}

Properties that always link to structured entities (100\% struct):
\texttt{itemReviewed}, \texttt{isPartOf}, \texttt{correction}, \texttt{video},
\texttt{sdPublisher}, \texttt{article}, \texttt{copyrightHolder},
\texttt{locationCreated}, \texttt{mentions}, \texttt{author}.

\subsection{Characteristic Property Sets}

\begin{table}[H]
\centering
\caption{Top 15 ClaimReview characteristic property sets (105 unique sets total).}
\label{tab:cr_charsets}
\begin{tabular}{rlL{9cm}}
\toprule
\textbf{\#} & \textbf{n (\%)} & \textbf{Property set} \\
\midrule
 1 & 22{,}896 (38.0\%) & \{none, only \texttt{rdf:type}\} \\
 2 & 16{,}737 (27.8\%) & \{author, claimReviewed, datePublished, itemReviewed, reviewRating, url\} \\
 3 &  3{,}885 (6.4\%)  & \{author, claimReviewed, dateModified, datePublished, image, itemReviewed, mainEntityOfPage, publisher, reviewRating, url\} \\
 4 &  3{,}335 (5.5\%)  & \{author, claimReviewed, itemReviewed, reviewRating, url\} \\
 5 &  2{,}329 (3.9\%)  & \{author, claimReviewed, datePublished, itemReviewed, name, reviewRating, url\} \\
 6 &  1{,}719 (2.9\%)  & \{author, claimReviewed, datePublished, description, itemReviewed, reviewRating, url\} \\
 7 &  1{,}091 (1.8\%)  & \{author, claimReviewed, datePublished, reviewRating, url\} \\
 8 &    963 (1.6\%)    & \{author, claimReviewed, datePublished, description, itemReviewed, name, reviewRating, url\} \\
 9 &    878 (1.5\%)    & \{author, claimReviewed, datePublished, itemReviewed, reviewRating, sdPublisher, url\} \\
10 &    626 (1.0\%)    & \{author, claimReviewed, dateModified, datePublished, description, headline, image, inLanguage, itemReviewed, keywords, mainEntityOfPage, name, publisher, reviewRating, url\} \\
11 &    488 (0.8\%)    & \{description, image, name, potentialAction, url\} \\
12 &    427 (0.7\%)    & \{author, claimReviewed, dateModified, datePublished, description, itemReviewed, reviewBody, reviewRating, url\} \\
13 &    386 (0.6\%)    & \{author, claimReviewed, dateModified, datePublished, description, itemReviewed, keywords, reviewBody, reviewRating, url\} \\
14 &    368 (0.6\%)    & \{author, claimReviewed, dateCreated, dateModified, datePublished, itemReviewed, reviewRating, url\} \\
15 &    340 (0.6\%)    & \{author, claimReviewed, datePublished, image, inLanguage, isPartOf, itemReviewed, mainEntityOfPage, publisher, reviewRating, url\} \\
\bottomrule
\end{tabular}
\end{table}

Set \#1 (38\%) is entirely the \texttt{vpravda.ru} cluster: entities carry only
a \texttt{ClaimReview} type declaration with no semantic properties and are
effectively noise. 

\subsection{Structure of \texttt{itemReviewed}}

\begin{table}[H]
\centering
\caption{Types of entities linked via \texttt{itemReviewed} (33{,}773 entities).}
\label{tab:ir_types}
\begin{tabular}{lrr}
\toprule
\textbf{Type} & \textbf{Count} & \textbf{\%} \\
\midrule
Claim        & 19{,}824 & 58.7\% \\
CreativeWork & 13{,}930 & 41.2\% \\
creativeWork (invalid capitalisation) &  17 & 0.1\% \\
NewsArticle  &       2  & 0.0\% \\
\bottomrule
\end{tabular}
\end{table}

\begin{table}[H]
\centering
\caption{Top 10 characteristic property sets for \texttt{itemReviewed} targets.}
\label{tab:ir_charsets}
\begin{tabular}{rlL{8.5cm}}
\toprule
\# & n (\%) & Property set \\
\midrule
 1 & 12{,}832 (38.0\%) & \{appearance, author, datePublished\} \\
 2 &  7{,}156 (21.2\%) & \{author, datePublished\} \\
 3 &  3{,}885 (11.5\%) & \{author, dateModified, datePublished, name\} \\
 4 &  2{,}532 (7.5\%)  & \{author, datePublished, url\} \\
 5 &  1{,}719 (5.1\%)  & \{author, datePublished, keywords\} \\
 6 &  1{,}007 (3.0\%)  & \{author, datePublished, image, name\} \\
 7 &    736 (2.2\%)    & \{appearance, author, datePublished, firstAppearance\} \\
 8 &    647 (1.9\%)    & \{author\} \\
 9 &    644 (1.9\%)    & \{datePublished\} \\
10 &    618 (1.8\%)    & \{author, dateModified, datePublished, image, mainEntityOfPage, name\} \\
\bottomrule
\end{tabular}
\end{table}


\subsection{Domain-Level Annotation Quality}
\label{sec:cr:domains}

\begin{table}[H]
\centering
\caption{Domain-level ClaimReview annotation quality (top 20 by entity count).
  \textbf{\%}: share of all 60{,}313 ClaimReview entities;
  \textbf{fill\%}: avg fraction of the 139 defined ClaimReview properties
  populated per entity (overall avg: 3.1\%).}
\label{tab:cr_domains}
\small
\begin{tabular}{L{3.2cm} R{1.2cm} R{1.0cm} R{1.2cm}}
\toprule
\textbf{Domain} & \textbf{n} & \textbf{\%} & \textbf{fill\%} \\
\midrule
vpravda.ru         & 22{,}891 & 38.0\% & 0.0\% \\
afp.com            &  2{,}918 &  4.8\% & 4.7\% \\
co.th              &  2{,}337 &  3.9\% & 4.3\% \\
boomlive.in        &  2{,}152 &  3.6\% & 7.2\% \\
leadstories.com    &  1{,}720 &  2.9\% & 5.0\% \\
org.pl             &  1{,}667 &  2.8\% & 4.3\% \\
healthday.com      &  1{,}330 &  2.2\% & 4.3\% \\
provereno.media    &  1{,}255 &  2.1\% & 4.3\% \\
types-cars.com     &  1{,}157 &  1.9\% & 4.3\% \\
vishvasnews.com    &  1{,}000 &  1.7\% & 4.3\% \\
misbar.com         &    963   &  1.6\% & 5.8\% \\
com.br             &    894   &  1.5\% & 4.7\% \\
fullfact.org       &    855   &  1.4\% & 6.8\% \\
politifact.com     &    787   &  1.3\% & 4.3\% \\
maldita.es         &    769   &  1.3\% & 3.6\% \\
tempo.co           &    629   &  1.0\% & 3.6\% \\
truthorfiction.com &    609   &  1.0\% & 3.6\% \\
correctiv.org      &    577   &  1.0\% & 4.3\% \\
newsmeter.in       &    567   &  0.9\% & 7.2\% \\
facta.news         &    567   &  0.9\% & 4.3\% \\
\bottomrule
\end{tabular}
\end{table}

The corpus-wide average fill rate is \textbf{3.1\%} across all 60{,}313
ClaimReview entities, driven almost entirely downward by \textbf{vpravda.ru}
(Volgograd regional newspaper, Russia), which contributes 38\% of all entities
with a fill rate of 0\%.
Excluding this cluster, established fact-checking organisations reach fill
rates of 3.6--7.2\%, with \texttt{boomlive.in} (7.2\%) and
\texttt{fullfact.org} (6.8\%) annotating most richly relative to the
defined vocabulary.


\subsection{Example Entities}

\paragraph{Depth-2 entity (canonical ClaimReview structure).}
\begin{quote}\small
\texttt{types: ClaimReview}\\
\texttt{graph: https://pagellapolitica.it/fact-checking/...}\\
\texttt{author} $\to$ \texttt{[Organization]\{name="Pagella Politica", url=...\}}\\
\texttt{claimReviewed} = ``Ecco i costi della regolarizzazione...''\\
\texttt{itemReviewed} $\to$ \texttt{[Claim]\{appearance=..., author=$\to$[Person], datePublished="2020-06-4"\}}\\
\texttt{reviewRating} $\to$ \texttt{[Rating]\{alternateName="Ni'"\}}\\
\texttt{url} = \texttt{<https://pagellapolitica.it/...>}
\end{quote}

\paragraph{Flat entity (depth~0, vpravda.ru noise).}
\begin{quote}\small
\texttt{types: ClaimReview}\\
\texttt{graph: https://vpravda.ru/1-maya/}\\
(no schema.org properties)
\end{quote}

\paragraph{Flat entity (depth~0, misapplied markup).}
\begin{quote}\small
\texttt{types: ClaimReview}\\
\texttt{graph: https://www.swico.ch/de/}\\
\texttt{claimReviewed} = ``Mitgliederfirmen vom Start-Up bis zum internationalen Konzern...''\\
\texttt{reviewRating} = \texttt{"700"@de} \quad \small(numeric membership count used as verdict)\\
\texttt{url} = \texttt{<https://www.swico.ch/de/...>}
\end{quote}

% =============================================================================
\section{Key Findings}
\label{sec:findings}
% =============================================================================

\begin{enumerate}[leftmargin=1.5em]

\item \textbf{High syntactic quality, low semantic richness.}
  Namespace correctness (99.9\%) and type validity (99.76\%) confirm that
  publishers generally deploy Schema.org identifiers correctly.
  Yet corpus-wide mean fill rate is just 7.4\% and mean AvP is 3.7
  (median 3.0): the macro-averaged property coverage across all 137 types
  is $5.75\%$ ($\sigma = 4.64\%$).  Publishers get the identifiers right
  but leave the vast majority of applicable vocabulary fields empty.

\item \textbf{Metadata, not content.}
  On article pages, structured data describes publication infrastructure
  (author, publisher, image, breadcrumbs) rather than semantic content.
  The properties most relevant to knowledge-graph grounding are almost
  entirely absent: \texttt{mentions} fills at effectively 0\% across all
  340{,}789 lenient article pages; \texttt{about} reaches only 19.3\%
  combined (39.6\% on \texttt{Article}, 1.3\% on \texttt{NewsArticle}),
  and 99.3\% of those targets are bare \texttt{Thing} nodes carrying only
  a free-text name string with no further structured properties;
  \texttt{citation} appears on 11.1\% of pages.
  These systematic gaps are the primary motivation for LLM-based enrichment.

\item \textbf{Claim entities are structurally hollow.}
  The \texttt{Claim} entity is the direct semantic payload of a
  fact-check---the structured representation of what claim was reviewed.
  Yet it has the lowest fill rate of any major type: 1.98\% (AvP 2.62 out
  of 132 defined properties).  Its dominant property set,
  \{appearance, author, datePublished\} (63.2\% of all \texttt{Claim}
  entities), records provenance but not the claim text itself; the claim
  statement is instead embedded in the parent \texttt{ClaimReview}'s
  \texttt{claimReviewed} string, which is itself absent from 40.3\% of
  all ClaimReview entities.

\item \textbf{ClaimReview quality is bimodal.}
  Established fact-checking organisations produce clean, deeply nested
  records: excluding the \texttt{vpravda.ru} noise cluster, the six core
  semantic properties (\texttt{author}, \texttt{claimReviewed},
  \texttt{datePublished}, \texttt{itemReviewed}, \texttt{reviewRating},
  \texttt{url}) are populated on approximately 96\% of remaining entities,
  and median nesting depth is 2.  In stark contrast, \texttt{vpravda.ru}
  inflates entity counts by 38\% while contributing zero semantic signal:
  every entity carries only an \texttt{rdf:type} declaration.  Even among
  the best-annotated organisations, fill rates peak at 7.2\%
  (\texttt{boomlive.in}) and 6.8\% (\texttt{fullfact.org}),
  indicating substantial room for automated enrichment.

\item \textbf{Range violations concentrate in two structural patterns.}
  Of the 11.0\% range-violation rate (6.05M violations from 55.1M checked
  assertions), 54\% are pixel integers in \texttt{height}/\texttt{width}
  where Schema.org expects a \texttt{QuantitativeValue} node, and most of
  the remainder are non-ISO date strings (e.g.\ \texttt{"Sep 14, 2024"},
  \texttt{"10/08/2023"}).  Neither is semantically critical for the
  journalism use case, but both complicate automated parsing.

\item \textbf{Verdict labels are uncontrolled.}
  \texttt{reviewRating.alternateName} values appear in multiple languages,
  scripts, and granularity levels with no standardised controlled vocabulary
  enforced by Schema.org.  Cross-organisation comparison or aggregation
  requires normalisation to a shared verdict taxonomy.

\item \textbf{Schema.org vocabulary vastly underutilised.}
  Only 41 of the 139 properties defined for \texttt{ClaimReview} (29.5\%)
  are ever used by any publisher.  Rich vocabulary items such as
  \texttt{NewsArticle.mentions}, \texttt{Article.about},
  \texttt{ClaimReview.associatedClaimReview}, and
  \texttt{Claim.appearance} are left entirely or almost entirely
  unexploited despite being directly relevant to fact-checking.

\item \textbf{Template-driven, publisher-uniform annotation.}
  Characteristic property sets are highly concentrated: the single most
  frequent set accounts for 28.4\% of \texttt{Article} entities, 22.8\%
  of \texttt{NewsArticle} entities, and 27.8\% of non-noise
  \texttt{ClaimReview} entities.  This uniformity reflects CMS-generated
  rather than editorially curated Schema.org output, meaning that
  completeness gaps are structural and cannot be closed by publisher
  incentives alone.

\end{enumerate}

% =============================================================================
\section{Implications for LLM-Based Triple Extraction}
\label{sec:implications}
% =============================================================================

The analyses above reveal a systematic gap between syntactically correct
Schema.org annotations and semantically useful structured data on news and
fact-checking pages.  This section translates each major finding into a
concrete requirement or opportunity for LLM-based enrichment.

\paragraph{Direct extraction targets on article pages.}
The three highest-priority missing properties are \texttt{mentions}
($<$0.1\% fill, which would provide named-entity links), \texttt{about}
(19.3\% combined fill, the primary topic entity), and \texttt{keywords}
(sparse and uncontrolled).  An LLM operating on \texttt{articleBody}
text --- present on 47.8\% of article pages --- can extract these
properties directly.  For the remaining 52.2\% of pages without a
structured body, the pipeline must first recover the article text from
raw HTML before extraction can proceed.

\paragraph{Entity typing and enrichment for \texttt{about} targets.}
Of the 65{,}146 Article and NewsArticle entities that do supply an
\texttt{about} link, 99.3\% point to bare \texttt{Thing} nodes carrying
only a free-text name string.  LLMs can replace these with properly
typed, property-populated entities (e.g.\ \texttt{Person},
\texttt{Organization}, \texttt{Place}) and augment them with
\texttt{sameAs} links to knowledge-base identifiers (Wikidata, DBpedia),
converting keyword strings into grounded knowledge-graph nodes.

\paragraph{Claim text and provenance completion.}
The claim statement (\texttt{claimReviewed}) is absent from 40.3\% of
\texttt{ClaimReview} entities; \texttt{Claim.appearance} --- linking to
where the claim first circulated --- is present in only 38\% of
\texttt{itemReviewed} targets.  LLMs can synthesise both from the
surrounding fact-check article text, filling in the two properties most
directly needed to reconstruct what claim was reviewed and where it
originated.

\paragraph{Verdict normalisation.}
The \texttt{reviewRating.alternateName} field is populated with
uncontrolled multilingual free-text labels.  LLM-based classification
can normalise these to a shared verdict taxonomy (e.g.\
True\,/\,Mostly True\,/\,Mostly False\,/\,False), enabling
cross-organisation comparison and aggregation that is currently
infeasible from the raw data.

\paragraph{Quality correction of systematic violations.}
The 11\% range-violation rate and the identified domain mismatches
(e.g.\ \texttt{MedicalProcedure}/\texttt{price} from a mis-typed entity,
\texttt{Article}/\texttt{primaryImageOfPage}) are structural defects an
LLM could detect and repair.  Correcting non-ISO date strings in
\texttt{datePublished} alone would eliminate 83{,}558 range violations
in the news-article subset, improving downstream parseability.

\paragraph{Priority vocabulary for Schema.org-aligned output.}
Since only 41 of 139 ClaimReview properties are ever used, LLM extraction
need not attempt full schema coverage.  A targeted property set ---
\texttt{claimReviewed}, \texttt{reviewRating}, \texttt{itemReviewed},
\texttt{author}, \texttt{datePublished}, \texttt{url},
\texttt{Claim.appearance}, \texttt{Claim.author}, and
\texttt{Article.mentions} --- covers the semantically critical vocabulary
while remaining tractable for prompt engineering and output parsing.

\paragraph{Existing annotations as evaluation baseline.}
The WDC data provides partial ground truth for a subset of properties.
Structured \texttt{ClaimReview} records from established organisations
(e.g.\ \texttt{boomlive.in}, \texttt{fullfact.org},
\texttt{politifact.com}) can serve as reference targets for evaluating
LLM extraction: a model given the raw fact-check article text should
reproduce the author, verdict, and \texttt{claimReviewed} string already
recorded in the markup.  The $\approx$96\% fill rate of core fields in
non-noise entities makes these examples reliable evaluation instances.

% =============================================================================
\section{End-to-End Summary}
\label{sec:summary}
% =============================================================================

Table~\ref{tab:e2e} collects the headline statistics from each stage of
the pipeline.

\begin{table}[H]
\centering
\caption{End-to-end pipeline summary.}
\label{tab:e2e}
\begin{tabular}{lr}
\toprule
\textbf{Stage / Metric} & \textbf{Value} \\
\midrule
\multicolumn{2}{l}{\textit{Target domain list}} \\
Total domains & 388 (390 PLDs) \\
Fact-check sources & Fact-Check Insights (top-15 orgs) \\
News sources & Media Cloud national high-volume \\
Manually added & 6 core domains \\
\midrule
\multicolumn{2}{l}{\textit{WDC chunk planning}} \\
Classes queried & 50 \\
Classes with matching domains & 37 \\
Total chunks (na\"ive plan) & 1{,}093 \\
Chunks after priority filter & 134 ($-$88\%) \\
\midrule
\multicolumn{2}{l}{\textit{Extraction}} \\
Total raw quads matched & 78{,}940{,}582 \\
Unique quads after deduplication & 69{,}621{,}246 \\
Uncompressed merged file & 14.4 GB \\
Distinct Schema.org types & 162 \\
Unique registered domains & 119 \\
\midrule
\multicolumn{2}{l}{\textit{Full-corpus quality (DuckDB, 18.6 GB)}} \\
Entities & 12{,}211{,}353 \\
Namespace correctness & 99.90\% \\
Type validity & 99.76\% \\
Property validity & 99.18\% \\
Range violation rate & 11.00\% \\
Mean AvP & 3.70 (median 3.0) \\
Mean fill rate & 7.4\% \\
Lenient article pages & 340{,}789 \\
\midrule
\multicolumn{2}{l}{\textit{ClaimReview (full WDC subset)}} \\
Entities & 60{,}313 \\
Domains & 318 \\
Namespace correctness & 100.00\% \\
Type validity & 100.00\% \\
Property validity & 99.74\% \\
Range violations & 2.27\% \\
Mean AvP & 4.24 (median 6.0) \\
Mean nesting depth & 1.54 (median 2.0) \\
Properties used / defined & 41 / 139 (29.5\%) \\
\bottomrule
\end{tabular}
\end{table}

\noindent
The pipeline begins with a curated list of 388 news and fact-checking domains
assembled from Fact-Check Insights and Media Cloud.  Using WDC's PLD lookup
tables, the na\"ive download plan of 1{,}093 chunks was reduced by 88\% to
134 chunks through a priority filter that preserves full coverage of core
domains while capping class-level downloads.  The resulting 69.6M unique
N-Quads spanning 119 registered domains were loaded into an 18.6 GB DuckDB
store, circumventing the out-of-memory ceiling that would otherwise block
in-process analysis of a 14.4 GB flat file.

The quality picture is one of syntactic competence paired with semantic
sparsity.  Namespace and type correctness exceed 99.7\%; property violations
are below 1\%.  Range violations reach 11\%, but 54\% of these collapse to
a single trivial pattern (pixel integers in \texttt{height}/\texttt{width}).
Yet the mean fill rate is just 7.4\% and mean AvP is 3.7: publishers emit
structurally valid Schema.org markup but populate only a small fraction of
the available vocabulary.  On the 340{,}789 article pages most relevant to
automated fact-checking, the properties with the highest semantic value ---
\texttt{mentions}, \texttt{about}, and \texttt{citation} --- fill at
$<$0.1\%, 19.3\%, and 11.1\% respectively, and nearly all \texttt{about}
targets are bare, untyped name strings.

The ClaimReview analysis reinforces this picture from the fact-checking side.
Established organisations produce clean, deeply nested records with
approximately 96\% coverage of the six core semantic properties, but a single
noise cluster (\texttt{vpravda.ru}) inflates entity counts by 38\% while
contributing zero semantic content.  Even the best-annotated organisations
fill only 7.2\% of the 139 defined ClaimReview properties, and the
\texttt{Claim} entity --- the structured representation of the reviewed claim
itself --- has an AvP of only 2.62 and a fill rate of 1.98\%.  Verdict labels
remain uncontrolled multilingual strings with no shared taxonomy.

Together, these results establish the baseline against which LLM-based triple
extraction should be measured.  The gap is systematic, not incidental: it
arises from CMS-templated annotation that is uniform within publishers but
shallow across the Schema.org vocabulary.  LLMs operating on article text can
in principle populate \texttt{mentions}, type and enrich \texttt{about}
targets, complete \texttt{claimReviewed} and \texttt{Claim.appearance}, and
normalise verdict labels --- precisely the properties that existing markup
leaves empty --- while the richer ClaimReview records from established
fact-checkers provide a reliable evaluation baseline for measuring extraction
quality.

\bibliographystyle{unsrtnat}
\bibliography{wdc_analysis_report}

\end{document}
